\documentclass[12pt,a4paper]{article}

% --- PACOTES FUNDAMENTAIS ---
\usepackage[utf8]{inputenc}
\usepackage[T1]{fontenc}
\usepackage[brazil]{babel}
\usepackage{geometry}
\geometry{a4paper, margin=2.5cm, top=3cm, bottom=2.5cm}

\usepackage{amsmath,amssymb}
\usepackage{booktabs}
\usepackage{tabularx}
\usepackage{graphicx}
\usepackage{xcolor}
\usepackage{listings}
\usepackage{hyperref}
\usepackage{fancyhdr}
\usepackage{float}
\usepackage{microtype}
\usepackage{tikz}
\usetikzlibrary{shapes.geometric, arrows.meta, positioning}

% --- CORES E CONFIGURAÇÕES VISUAIS ---
\definecolor{primary}{RGB}{16, 52, 166}
\definecolor{secondary}{RGB}{44, 62, 80}
\definecolor{codebg}{RGB}{248, 249, 250}
\definecolor{codeborder}{RGB}{220, 224, 230}
\definecolor{codegreen}{RGB}{46, 125, 50}
\definecolor{codegray}{RGB}{120, 144, 156}
\definecolor{codepurple}{RGB}{123, 31, 162}

\hypersetup{
    colorlinks=true,
    linkcolor=primary,
    filecolor=primary,
    urlcolor=primary,
    citecolor=secondary
}

% --- CONFIGURAÇÃO DE LISTINGS (CÓDIGO) ---
\lstset{
    backgroundcolor=\color{codebg},
    basicstyle=\ttfamily\footnotesize,
    breaklines=true,
    captionpos=b,
    commentstyle=\color{codegray}\itshape,
    keywordstyle=\color{primary}\bfseries,
    stringstyle=\color{codegreen},
    numberstyle=\tiny\color{codegray},
    numbers=left,
    numbersep=8pt,
    frame=single,
    rulecolor=\color{codeborder},
    xleftmargin=15pt,
    framexleftmargin=15pt,
    tabsize=4,
    showstringspaces=false
}

% --- CABEÇALHO E RODAPÉ ---
\pagestyle{fancy}
\fancyhf{}
\lhead{\footnotesize \textbf{Sistemas Distribuídos} — Relatório de Implementação}
\rhead{\footnotesize Marcos 1 e 2}
\cfoot{\thepage}
\renewcommand{\headrulewidth}{0.4pt}

% --- INÍCIO DO DOCUMENTO ---
\begin{document}

\begin{titlepage}
    \centering
    \vspace*{1.5cm}
    
    {\scshape\LARGE Universidade Federal / Instituição de Ensino Superior \par}
    \vspace{0.5cm}
    {\scshape\Large Disciplina de Sistemas Distribuídos \par}
    \vspace{2cm}
    
    \rule{\linewidth}{1.5pt}\vspace{0.4cm}
    {\huge\bfseries Sistema Distribuído para Auxílio de Dados Hospitalares\par}
    \vspace{0.3cm}
    {\Large\itshape Relatório Técnico e Descritivo de Implementação: Marcos 1 e 2\par}
    \rule{\linewidth}{1.5pt}
    
    \vspace{2.5cm}
    
    \begin{minipage}{0.8\textwidth}
        \centering
        \textbf{Autor:} Eduardo \\
        \textbf{Contexto:} Projeto Acadêmico Prático \\
        \textbf{Foco:} Telemetria Hospitalar, APIs Assíncronas, Persistência Local e Balanceamento de Carga Round Robin
    \end{minipage}
    
    \vfill
    
    {\large \today \par}
\end{titlepage}

\newpage
\tableofcontents
\newpage

% =========================================================================
\section{Introdução e Visão Geral}
% =========================================================================

O monitoramento contínuo de pacientes em ambientes hospitalares gera fluxos ininterruptos de telemetria crítica (sinais vitais, alertas fisiológicos e ocorrências de emergência). Em arquiteturas computacionais centralizadas, esse volume de requisições pode sobrecarregar servidores de borda, introduzir latência no processamento e criar pontos únicos de falha (\textit{Single Points of Failure}).

O presente projeto tem por finalidade o projeto e a implementação de um \textbf{Sistema Distribuído para Auxílio de Dados Hospitalares}, concebido com propósitos estritamente acadêmicos e didáticos. A proposta adota uma abordagem incremental estruturada em marcos evolutivos. Este documento detalha a concepção, formalização e validação dos dois primeiros marcos:

\begin{itemize}
    \item \textbf{Marco 1 — Simulação e Servidor Base:} Ingestão progressiva de telemetria real proveniente de um conjunto de dados hospitalar (\textit{dataset} Kaggle), validação estrita de esquemas em camada de API RESTful (FastAPI), conversão estruturada de tipos e persistência relacional local (SQLite).
    \item \textbf{Marco 2 — Balanceamento de Carga e Múltiplas Instâncias:} Criação de três instâncias concorrentes do servidor e de um Balanceador de Carga em camada de aplicação implementando o algoritmo \textit{Round Robin}, com transparência de acesso e identificação inequívoca de nós processadores.
\end{itemize}

O projeto foi integralmente construído em linguagem Python (versão 3.10+), priorizando clareza conceitual, desacoplamento e portabilidade.

% =========================================================================
\section{Análise e Modelagem do Dataset Hospitalar}
% =========================================================================

\subsection{Origem e Caracterização do Arquivo}
A fonte de dados empregada é o arquivo CSV obtido via Kaggle:
\begin{center}
\texttt{Synthetic\_patient-HealthCare-Monitoring\_dataset.csv}
\end{center}
Diferente de geradores sintéticos aleatórios gerados em tempo de execução, este arquivo atua como repositório de telemetria pré-coletada.

Uma análise exploratória minuciosa do arquivo revelou:
\begin{itemize}
    \item \textbf{Volume Total:} 60.000 registros clínicos tabulados (60.001 linhas contando o cabeçalho).
    \item \textbf{Completude dos Dados:} Não foram identificados registros nulos ou incompletos (0\% de valores ausentes em todas as colunas).
    \item \textbf{Ausência de Carimbo de Tempo:} O conjunto original não possui coluna temporal (\textit{timestamp}), exigindo que a camada de borda do servidor gere um carimbo no momento exato do recebimento da mensagem.
\end{itemize}

\subsection{Camada de Mapeamento (Schema Mapping)}
Para adequar a terminologia do arquivo (em língua inglesa e com caracteres especiais) às boas práticas de Engenharia de Software e APIs RESTful em português, foi implementado um mapeamento bidirecional no módulo \texttt{server/models.py}.

\begin{table}[H]
\centering
\footnotesize
\caption{Mapeamento de Atributos do CSV Kaggle para a API e Banco de Dados}
\label{tab:schema_map}
\begin{tabularx}{\textwidth}{p{4.2cm} p{4.2cm} p{1.2cm} X}
\toprule
\textbf{Coluna Original (CSV)} & \textbf{Campo na API/DB} & \textbf{Tipo} & \textbf{Descrição Semântica} \\
\midrule
\texttt{Patient Number} & \texttt{paciente\_id} & \texttt{int} & Identificador do paciente \\
\texttt{Heart Rate (bpm)} & \texttt{frequencia\_cardiaca} & \texttt{int} & Frequência cardíaca (bpm) \\
\texttt{SpO2 Level (\%)} & \texttt{spo2} & \texttt{int} & Saturação de oxigênio \\
\texttt{Systolic Blood Pressure} & \texttt{pressao\_sistolica} & \texttt{int} & Pressão sistólica (mmHg) \\
\texttt{Diastolic Blood Pressure} & \texttt{pressao\_diastolica} & \texttt{int} & Pressão diastólica (mmHg) \\
\texttt{Body Temperature (°C)} & \texttt{temperatura} & \texttt{float} & Temperatura corpórea (°C) \\
\texttt{Fall Detection} & \texttt{deteccao\_queda} & \texttt{str} & Indicador de queda (Yes/No) \\
\texttt{Predicted Disease} & \texttt{doenca\_prevista} & \texttt{str} & Diagnóstico clínico previsto \\
\texttt{Data Accuracy (\%)} & \texttt{precisao\_dados} & \texttt{int} & Confiabilidade do sensor (\%) \\
\texttt{Heart Rate Alert} & \texttt{alerta\_freq\_cardiaca} & \texttt{str} & Alerta de FC (NORMAL/ABNORMAL) \\
\texttt{SpO2 Level Alert} & \texttt{alerta\_spo2} & \texttt{str} & Alerta de SpO2 \\
\texttt{Blood Pressure Alert} & \texttt{alerta\_pressao} & \texttt{str} & Alerta de pressão arterial \\
\texttt{Temperature Alert} & \texttt{alerta\_temperatura} & \texttt{str} & Alerta de temperatura \\
\bottomrule
\end{tabularx}
\end{table}

% =========================================================================
\section{Marco 1 — Simulação e Servidor Base}
% =========================================================================

\subsection{Arquitetura e Fluxo de Informação}
O Marco 1 estabelece o pipeline elementar de telemetria hospitalar unidirecional. A Figura~\ref{fig:fluxo_marco1} ilustra os componentes e o ciclo de vida do dado.

\begin{figure}[H]
\centering
\resizebox{\textwidth}{!}{%
\begin{tikzpicture}[
    block/.style={rectangle, draw=primary, fill=codebg, thick, rounded corners, minimum height=1.1cm, minimum width=2.8cm, align=center},
    arrow/.style={-Latex, thick, primary}
]
    \node[block] (csv) {\textbf{Arquivo CSV}\\[2pt]\footnotesize 60.000 Registros};
    \node[block, right=2.2cm of csv] (sender) {\textbf{Leitor / Produtor}\\[2pt]\footnotesize \texttt{csv\_sender.py}};
    \node[block, right=3.2cm of sender] (server) {\textbf{Servidor Base}\\[2pt]\footnotesize FastAPI (:8001)};
    \node[block, right=2.4cm of server] (db) {\textbf{SQLite Local}\\[2pt]\footnotesize \texttt{server1.db}};

    \draw[arrow] (csv) -- node[above=2pt, font=\footnotesize\bfseries, color=primary]{Leitura} (sender);
    \draw[arrow] (sender) -- node[above=2pt, font=\footnotesize\bfseries, color=primary]{HTTP POST} node[below=2pt, font=\footnotesize\ttfamily, color=primary]{/dados} (server);
    \draw[arrow] (server) -- node[above=2pt, font=\footnotesize\bfseries, color=primary]{INSERT} (db);
\end{tikzpicture}%
}
\caption{Fluxo de telemetria do Marco 1: do arquivo tabular à persistência local.}
\label{fig:fluxo_marco1}
\end{figure}

\subsection{O Produtor de Telemetria (\texttt{client/csv\_sender.py})}
O script atua como um emulador de dispositivos biomédicos de leito. Suas responsabilidades incluem:
\begin{enumerate}
    \item Abertura segura do arquivo CSV via \texttt{csv.DictReader};
    \item Conversão atômica de cada linha através da função \texttt{converter\_registro\_csv}, efetuando o \textit{casting} de strings para inteiros e números de ponto flutuante;
    \item Envio progressivo via requisição síncrona HTTP \texttt{POST} com corpo em formato JSON para o endpoint de destino;
    \item Temporização configurável entre eventos (\texttt{time.sleep(intervalo)}), viabilizando simulações aceleradas ou em tempo real;
    \item Tratamento resiliente de falhas de conexão (\texttt{ConnectionError} e \texttt{Timeout}), garantindo que falhas momentâneas de rede não interrompam o processo.
\end{enumerate}

\subsection{O Servidor Base (\texttt{server/app.py})}
Construído sobre o framework assíncrono \textbf{FastAPI}, o servidor expõe três interfaces essenciais:
\begin{itemize}
    \item \texttt{POST /dados}: Recebe a carga útil validada pelo modelo \texttt{DadosPaciente} (Pydantic). O servidor injeta o \texttt{timestamp} da máquina e seu identificador (\texttt{server\_id}), grava no banco e retorna HTTP 200 com a entidade persistida.
    \item \texttt{GET /dados}: Retorna registros armazenados, oferecendo suporte nativo à paginação através dos parâmetros de consulta \texttt{limit} e \texttt{offset} (essencial para evitar estouro de memória em datasets com dezenas de milhares de linhas).
    \item \texttt{GET /dados/\{paciente\_id\}}: Filtra o histórico longitudinal de um paciente específico, retornando HTTP 404 caso o identificador não tenha medições registradas.
\end{itemize}

\subsection{Camada de Persistência (\texttt{server/database.py})}
Utiliza o mecanismo \textbf{SQLite3}. Cada instância do servidor gerencia seu próprio arquivo de banco de dados, inicializando a tabela \texttt{dados\_pacientes} na inicialização:

\begin{lstlisting}[language=SQL, caption={Estrutura da Tabela no SQLite}]
CREATE TABLE IF NOT EXISTS dados_pacientes (
    id INTEGER PRIMARY KEY AUTOINCREMENT,
    paciente_id INTEGER NOT NULL,
    frequencia_cardiaca INTEGER NOT NULL,
    spo2 INTEGER NOT NULL,
    pressao_sistolica INTEGER NOT NULL,
    pressao_diastolica INTEGER NOT NULL,
    temperatura REAL NOT NULL,
    deteccao_queda TEXT NOT NULL,
    doenca_prevista TEXT NOT NULL,
    precisao_dados INTEGER NOT NULL,
    alerta_freq_cardiaca TEXT NOT NULL,
    alerta_spo2 TEXT NOT NULL,
    alerta_pressao TEXT NOT NULL,
    alerta_temperatura TEXT NOT NULL,
    timestamp TEXT NOT NULL,
    server_id TEXT NOT NULL
);
\end{lstlisting}

% =========================================================================
\section{Marco 2 — Balanceamento de Carga e Múltiplas Instâncias}
% =========================================================================

\subsection{Arquitetura Distribuída Proposta}
No Marco 2, a arquitetura é escalada horizontalmente. Introduz-se um nó intermediário (o Balanceador de Carga) na porta 8000, responsável por desacoplar o cliente produtor dos servidores de aplicação (instanciados nas portas 8001, 8002 e 8003).

\begin{figure}[H]
\centering
\resizebox{0.95\textwidth}{!}{%
\begin{tikzpicture}[
    node distance=1.2cm and 2.0cm,
    block/.style={rectangle, draw=primary, fill=codebg, thick, rounded corners, minimum height=1.1cm, minimum width=2.9cm, align=center},
    servernode/.style={rectangle, draw=secondary, fill=codebg, thick, rounded corners, minimum height=0.95cm, minimum width=2.6cm, align=center},
    dbnode/.style={cylinder, draw=primary, shape border rotate=90, aspect=0.25, fill=codebg, thick, minimum height=0.85cm, minimum width=2cm, align=center},
    arrow/.style={-Latex, thick, primary}
]
    \node[block] (sender) {\textbf{Leitor de Dados}\\[2pt]\footnotesize \texttt{csv\_sender.py}};
    \node[block, right=3.2cm of sender] (lb) {\textbf{Balanceador}\\[2pt]\footnotesize Porta 8000};
    
    \node[servernode, above right=0.4cm and 2.0cm of lb] (s1) {\textbf{Server 1}\\[1pt]\footnotesize :8001};
    \node[servernode, right=2.0cm of lb] (s2) {\textbf{Server 2}\\[1pt]\footnotesize :8002};
    \node[servernode, below right=0.4cm and 2.0cm of lb] (s3) {\textbf{Server 3}\\[1pt]\footnotesize :8003};

    \node[dbnode, right=1.2cm of s1] (db1) {\footnotesize \texttt{server1.db}};
    \node[dbnode, right=1.2cm of s2] (db2) {\footnotesize \texttt{server2.db}};
    \node[dbnode, right=1.2cm of s3] (db3) {\footnotesize \texttt{server3.db}};

    \draw[arrow] (sender) -- node[above=2pt, font=\footnotesize\bfseries, color=primary]{HTTP POST} (lb);
    \draw[arrow] (lb) -- (s1);
    \draw[arrow] (lb) -- (s2);
    \draw[arrow] (lb) -- (s3);

    \draw[arrow] (s1) -- (db1);
    \draw[arrow] (s2) -- (db2);
    \draw[arrow] (s3) -- (db3);
\end{tikzpicture}%
}
\caption{Topologia física e lógica do Marco 2: balanceamento em anel uniforme.}
\label{fig:topologia_marco2}
\end{figure}

\subsection{Algoritmo Round Robin (Formulação e Implementação)}
O algoritmo \textbf{Round Robin} (revezamento circular) distribui as tarefas sequencialmente ao longo de uma lista ordenada de nós disponíveis. Dado um conjunto estático de $N$ servidores $\mathcal{S} = \{s_0, s_1, \dots, s_{N-1}\}$ e um contador de requisições acumuladas $r \in \mathbb{N}$ (iniciado em 0):

\begin{equation}
    \text{Servidor Alvo}(r) = s_{r \pmod N}
\end{equation}

Esta regra garante distribuição perfeitamente equitativa da carga útil de requisições:
\begin{itemize}
    \item Requisição 1 ($r=0$): Servidor $0$ (\texttt{server-01} na porta 8001)
    \item Requisição 2 ($r=1$): Servidor $1$ (\texttt{server-02} na porta 8002)
    \item Requisição 3 ($r=2$): Servidor $2$ (\texttt{server-03} na porta 8003)
    \item Requisição 4 ($r=3$): Servidor $0$ (\texttt{server-01} na porta 8001)
\end{itemize}

O balanceador foi implementado através da classe \texttt{RoundRobinBalancer}, no módulo \texttt{load\_balancer/balancer.py}. Para o encaminhamento assíncrono das requisições, foi empregada a biblioteca \texttt{httpx.AsyncClient}.

\subsection{Identificação e Transparência}
Cada instância do servidor executa o mesmo código-fonte, porém é parametrizada via linha de comando com identidade e banco próprios:
\begin{itemize}
    \item \textbf{Instância 1:} Porta \texttt{8001}, ID \texttt{"server-01"}, banco \texttt{databases/server1.db}
    \item \textbf{Instância 2:} Porta \texttt{8002}, ID \texttt{"server-02"}, banco \texttt{databases/server2.db}
    \item \textbf{Instância 3:} Porta \texttt{8003}, ID \texttt{"server-03"}, banco \texttt{databases/server3.db}
\end{itemize}

Quando o balanceador despacha a requisição e recebe a resposta do nó selecionado, a resposta HTTP 200 propagada de volta ao cliente contém explicitamente o campo:
\begin{center}
\texttt{"server\_id": "server-01"}
\end{center}
Isso permite auditoria em tempo real pelo cliente da distribuição de carga sem quebrar a \textit{transparência de localização} (o cliente continua se comunicando unicamente com o endereço \texttt{http://localhost:8000}).

\subsection{Bancos de Dados Independentes no Marco 2}
Conforme estipulado no planejamento pedagógico do projeto, \textbf{nenhuma replicação ou sincronização entre bancos foi implementada neste estágio}. 

Cada servidor alimenta estritamente seu próprio banco SQLite local (\texttt{server1.db}, \texttt{server2.db} ou \texttt{server3.db}). Se 12 registros forem enviados pelo balanceador, cada um dos 3 bancos registrará precisamente 4 tuplas. A persistência distribuída consistente e a replicação de réplicas ativas/passivas constituem o escopo do futuro Marco 3.

% =========================================================================
\section{Decisões de Engenharia e Desafios Superados}
% =========================================================================

Durante o desenvolvimento dos Marcos 1 e 2, algumas decisões de engenharia de software foram determinantes para a robustez da solução:

\begin{enumerate}
    \item \textbf{Substituição de Manipuladores Depreciados do FastAPI:} Versões recentes do FastAPI descontinuaram o padrão decorador \texttt{@app.on\_event("startup")}. Adotou-se o gerenciador de contexto de ciclo de vida assíncrono (\texttt{lifespan}), combinado com a invocação antecipada e idempotente de \texttt{init\_db()}, permitindo que os testes com \texttt{TestClient} operassem sem falhas de tabela inexistente.
    \item \textbf{Preservação da Integridade do Dataset Original:} O arquivo CSV de 4,4 MB não foi duplicado nem modificado. Criou-se um \textit{symbolic link} (\texttt{data/dataset.csv}) referenciando o original, garantindo reproducibilidade sem redundância de armazenamento.
    \item \textbf{Isolamento de Concorrência nos Testes do Marco 2:} Em testes automatizados envolvendo servidores reais e balanceadores em segundo plano, a reutilização rápida de portas TCP pode causar o erro \texttt{Errno 98: Address already in use}. A solução adotada consistiu em empregar \textit{fixtures} de escopo de classe (\texttt{scope="class"}) no \texttt{pytest} com portas dinâmicas na faixa 19000, eliminando condições de corrida e travamentos em sockets.
\end{enumerate}

% =========================================================================
\section{Validação e Resultados Experimentais}
% =========================================================================

O sistema foi rigorosamente validado através de 26 testes automatizados construídos com \texttt{pytest}, subdivididos em duas suítes dedicadas.

\begin{table}[H]
\centering
\small
\caption{Sumário da Suíte de Testes Automatizados}
\label{tab:testes}
\begin{tabularx}{\textwidth}{l c X}
\toprule
\textbf{Módulo de Teste} & \textbf{Qtd.} & \textbf{Objetivo da Validação} \\
\midrule
\texttt{tests/test\_marco1.py} & 17 & Existência do CSV, presença de colunas, parsing sem perda, validação Pydantic, inserção no SQLite, queries com paginação (\texttt{limit/offset}) e busca por paciente. \\
\texttt{tests/test\_marco2.py} & 9 & Aritmética modular do Round Robin, inicialização simultânea de 3 servidores e do balanceador, roteamento HTTP transparente e partição homogênea nos bancos. \\
\midrule
\textbf{Total de Testes Aprovados} & \textbf{26/26} & \textbf{Taxa de Sucesso: 100\% (tempo de execução: 0,92s)} \\
\bottomrule
\end{tabularx}
\end{table}

% =========================================================================
\section{Roteiro Prático de Demonstração e Verificação}
\label{sec:roteiro}
% =========================================================================

Para apresentar o projeto ao professor ou avaliadores e comprovar inequivocamente que os Marcos 1 e 2 estão plenamente funcionais, deve-se seguir o procedimento estruturado a seguir.

\subsection{Etapa Prévia: Ambiente e Testes Rápidos}
No terminal principal, ative o ambiente virtual e execute a suíte de testes:
\begin{lstlisting}[language=bash]
cd "Projeto SD"
source venv/bin/activate
python -m pytest tests/ -v
\end{lstlisting}
\textbf{O que mostrar:} A saída verde confirmando que todos os 26 testes passaram com sucesso em menos de 1 segundo.

\subsection{Demonstração Prática do Marco 1 (Servidor Base)}
Para demonstrar o funcionamento elementar do Marco 1:
\begin{enumerate}
    \item \textbf{Terminal 1 (Servidor):}
    \begin{lstlisting}[language=bash]
    python server/app.py --port 8001 --server-id server-01 --db-path databases/server1.db
    \end{lstlisting}
    \item \textbf{Terminal 2 (Envio Direto):}
    \begin{lstlisting}[language=bash]
    python client/csv_sender.py --url http://localhost:8001/dados --intervalo 0.3 --limite 5
    \end{lstlisting}
    \item \textbf{O que explicar ao avaliador:}
    \begin{itemize}
        \item No Terminal 2, aponte que os registros 1 a 5 foram lidos do CSV real e enviados via HTTP POST.
        \item No Terminal 1, mostre os logs registrando o recebimento dos sinais vitais.
        \item Abra o navegador ou utilize o \texttt{curl} para demonstrar a consulta dos dados armazenados:
        \begin{lstlisting}[language=bash]
        curl -s http://localhost:8001/dados | python -m json.tool
        \end{lstlisting}
        \item Consulte o banco SQLite diretamente pelo terminal para provar a persistência:
        \begin{lstlisting}[language=bash]
        sqlite3 databases/server1.db "SELECT id, paciente_id, frequencia_cardiaca, timestamp, server_id FROM dados_pacientes;"
        \end{lstlisting}
    \end{itemize}
\end{enumerate}

\subsection{Demonstração Prática do Marco 2 (Balanceamento de Carga)}
Esta é a comprovação máxima do sistema distribuído. Feche a execução do teste anterior e prepare os 5 terminais:

\begin{itemize}
    \item \textbf{Terminal 1 (Server 1):}
    \begin{lstlisting}[language=bash]
    python server/app.py --port 8001 --server-id server-01 --db-path databases/server1.db
    \end{lstlisting}
    \item \textbf{Terminal 2 (Server 2):}
    \begin{lstlisting}[language=bash]
    python server/app.py --port 8002 --server-id server-02 --db-path databases/server2.db
    \end{lstlisting}
    \item \textbf{Terminal 3 (Server 3):}
    \begin{lstlisting}[language=bash]
    python server/app.py --port 8003 --server-id server-03 --db-path databases/server3.db
    \end{lstlisting}
    \item \textbf{Terminal 4 (Balanceador):}
    \begin{lstlisting}[language=bash]
    python load_balancer/balancer.py
    \end{lstlisting}
    \item \textbf{Terminal 5 (Cliente Produtor):}
    Envie um lote controlado de 12 requisições através do balanceador (porta 8000):
    \begin{lstlisting}[language=bash]
    python client/csv_sender.py --url http://localhost:8000/dados --intervalo 0.4 --limite 12
    \end{lstlisting}
\end{itemize}

\subsection{Evidências de Funcionamento que Devem Ser Apontadas}
\begin{enumerate}
    \item \textbf{Alternância Visual no Balanceador (Terminal 4):}
    O balanceador exibirá explicitamente o revezamento cíclico das 12 requisições:
\begin{lstlisting}[basicstyle=\ttfamily\scriptsize]
[BALANCEADOR] Requisicao 1 -> http://localhost:8001
[BALANCEADOR] Requisicao 2 -> http://localhost:8002
[BALANCEADOR] Requisicao 3 -> http://localhost:8003
[BALANCEADOR] Requisicao 4 -> http://localhost:8001
[BALANCEADOR] Requisicao 5 -> http://localhost:8002
...
\end{lstlisting}
    \item \textbf{Reconhecimento pelo Cliente (Terminal 5):}
    O leitor confirmará o processamento indicando qual servidor respondeu:
\begin{lstlisting}[basicstyle=\ttfamily\scriptsize]
[    1] Paciente     1 -> server-01 | FC=98  SpO2=96 PA=120/86 Temp=38.1
[    2] Paciente     2 -> server-02 | FC=105 SpO2=97 PA=177/104 Temp=37.6
[    3] Paciente     3 -> server-03 | FC=90  SpO2=85 PA=139/57 Temp=37.0
[    4] Paciente     4 -> server-01 | FC=102 SpO2=87 PA=101/77 Temp=36.4
\end{lstlisting}
    \item \textbf{Divisão Exata nos Três Bancos SQLite:}
    Abra um terminal adicional e execute a contagem em cada banco:
    \begin{lstlisting}[language=bash]
    echo "--- Server 1 (8001) ---"
    sqlite3 databases/server1.db "SELECT COUNT(*) FROM dados_pacientes;"
    echo "--- Server 2 (8002) ---"
    sqlite3 databases/server2.db "SELECT COUNT(*) FROM dados_pacientes;"
    echo "--- Server 3 (8003) ---"
    sqlite3 databases/server3.db "SELECT COUNT(*) FROM dados_pacientes;"
    \end{lstlisting}
    \textbf{Resultado Incontestável:} Cada um dos três bancos exibirá exatamente \textbf{4 registros} (totalizando os 12 enviados). Isso prova que a carga foi perfeitamente distribuída e que as instâncias operam com persistência isolada.
\end{enumerate}

% =========================================================================
\section{Conclusão}
% =========================================================================

Os objetivos previstos para os Marcos 1 e 2 foram plenamente alcançados. A solução atende a todos os requisitos didáticos de Sistemas Distribuídos:
\begin{itemize}
    \item Demonstração de heterogeneidade e acoplamento fraco mediante o protocolo HTTP e JSON;
    \item Preservação de dados reais de telemetria sem artifícios sintéticos desnecessários;
    \item Implementação de balanceamento de carga Round Robin com transparência de acesso e tolerância a falhas na camada cliente;
    \item Estruturação modular que viabiliza a evolução suave para o Marco 3 (sincronização e replicação de dados distribuídos).
\end{itemize}

\end{document}
